%%%%%%%%%%%%%%%%
%%% exod 4 %%%
%%%%%%%%%%%%%%%%

\Chapter{4}

\Verse{1}
{
	\hE{וַיַּעַן מֹשֶׁה וַיֹּאמֶר וְהֵן לֹא יַאֲמִינוּ לִי וְלֹא יִשְׁמְעוּ בְּקֹלִי כִּי יֹאמְרוּ לֹא נִרְאָה אֵלֶיךָ יְהֹוָה׃}
}
{
	\eE{
		And Moses answered, and he said, “And here, they won’t
		believe me and won’t listen to my voice, because they’ll
		say, ‘YHWH hasn’t appeared to you!’”
	}
}
{
%	\footnote{
%		they won’t listen to my voice. Moses’ third response is surprising
%		because he appears to doubt, even to contradict, what God has already
%		told him explicitly. YHWH instructs him to gather the elders of Israel
%		and tell them that their ancestors’ God has appeared to him and is going
%		to take them out of Egypt, and YHWH informs him directly that “They’ll
%		listen to your voice” (3:18). God then goes on for five verses with
%		details of what Moses should say to the king of Egypt, what the king’s
%		response will be, what YHWH will do in return, and how Israel will leave
%		and despoil Egypt. But Moses, appearing to have missed the content of
%		these last five verses, questions the earlier point, the last point that
%		directly referred to him himself. He says, “And here, they won’t believe
%		me and won’t listen to my voice” (4:1). This can also be understood to
%		mean “And what if they will not believe me and will not listen to my
%		voice.” Even on this understanding, it is remarkable. Either way, it is
%		depicting Moses as questioning what God has just told him. When God
%		Himself tells a human, “They’ll listen to you,” we do not normally
%		expect the human to say, “But what if they don’t?” Again YHWH does not
%		answer the question but instead responds directly to the problem at
%		hand. That is, instead of telling Moses what to do if they do not
%		listen, God gives Moses three miraculous signs, which will guarantee
%		that they will listen.
%	}
}

\Verse{2}
{
	\hE{וַיֹּאמֶר אֵלָיו יְהֹוָה (מזה) [מַה זֶּה] בְיָדֶךָ וַיֹּאמֶר מַטֶּה׃}
}
{
	\eE{
		And YHWH said to him, “What’s this in your hand?” And he
		said, “A staff.”
	}
}
{
	Qere Ketiv.
}

\Verse{3}
{
	\hE{וַיֹּאמֶר הַשְׁלִיכֵהוּ אַרְצָה וַיַּשְׁלִכֵהוּ אַרְצָה וַיְהִי לְנָחָשׁ וַיָּנס מֹשֶׁה מִפָּנָיו׃}
}
{
	\eE{
		And He said, “Throw it to the ground.” And he threw it to
		the ground. And it became a snake! And Moses fled from it.
	}
}
{
}

\Verse{4}
{
	\hE{וַיֹּאמֶר יְהֹוָה אֶל מֹשֶׁה שְׁלַח יָדְךָ וֶאֱחֹז בִּזְנָבוֹ וַיִּשְׁלַח יָדוֹ וַיַּחֲזֶק בּוֹ וַיְהִי לְמַטֶּה בְּכַפּוֹ׃}
}
{
	\eE{
		And YHWH said to Moses, “Put out your hand and take hold
		of its tail.” And he put out his hand and held onto it,
		and it became a staff in his hand.
	}
}
{
}

\Verse{5}
{
	\hE{לְמַעַן יַאֲמִינוּ כִּי נִרְאָה אֵלֶיךָ יְהֹוָה אֱלֹהֵי אֲבֹתָם אֱלֹהֵי אַבְרָהָם אֱלֹהֵי יִצְחָק וֵאלֹהֵי יַעֲקֹב׃}
}
{
	\eE{
		“So that they will believe that YHWH, their fathers’ God,
		Abraham’s God, Isaac’s God, and Jacob’s God, has appeared
		to you!”
	}
}
{
}

\Verse{6}
{
	\hE{וַיֹּאמֶר יְהֹוָה לוֹ עוֹד הָבֵא נָא יָדְךָ בְּחֵיקֶךָ וַיָּבֵא יָדוֹ בְּחֵיקוֹ וַיּוֹצִאָהּ וְהִנֵּה יָדוֹ מְצֹרַעַת כַּשָּׁלֶג׃}
}
{
	\eE{
		And YHWH said to him further, “Bring your hand into your
		bosom.” And he brought his hand into his bosom. And he
		brought it out; and, here, his hand was leprous like snow!
	}
}
{
}

\Verse{7}
{
	\hE{וַיֹּאמֶר הָשֵׁב יָדְךָ אֶל חֵיקֶךָ וַיָּשֶׁב יָדוֹ אֶל חֵיקוֹ וַיּוֹצִאָהּ מֵחֵיקוֹ וְהִנֵּה שָׁבָה כִּבְשָׂרוֹ׃}
}
{
	\eE{
		And He said, “Put your hand back to your bosom.” And he
		put his hand back to his bosom. And he brought it out from
		his bosom; and, here, it had gone back like its flesh.
	}
}
{
}

\Verse{8}
{
	\hE{וְהָיָה אִם לֹא יַאֲמִינוּ לָךְ וְלֹא יִשְׁמְעוּ לְקֹל הָאֹת הָרִאשׁוֹן וְהֶאֱמִינוּ לְקֹל הָאֹת הָאַחֲרוֹן׃}
}
{
	\eE{
		“And it will be, if they won’t believe you and won’t
		listen to the voice of the first sign, then they’ll
		believe the voice of the latter one.
	}
}
{
}

\Verse{9}
{
	\hE{וְהָיָה אִם לֹא יַאֲמִינוּ גַּם לִשְׁנֵי הָאֹתוֹת הָאֵלֶּה וְלֹא יִשְׁמְעוּן לְקֹלֶךָ וְלָקַחְתָּ מִמֵּימֵי הַיְאֹר וְשָׁפַכְתָּ הַיַּבָּשָׁה וְהָיוּ הַמַּיִם אֲשֶׁר תִּקַּח מִן הַיְאֹר וְהָיוּ לְדָם בַּיַּבָּשֶׁת׃}
}
{
	\eE{
		And it will be, if they also won’t believe in these two
		signs and won’t listen to your voice, then you’ll take
		some of the water of the Nile and spill it on the dry
		ground. And it will be water that you’ll take from the
		Nile, and it will become blood on the dry ground.”
	}
}
{
%	\footnote{
%		two signs. The two signs are that the staff becomes a snake and his hand
%		becomes leprous as snow. Why these two things? They foreshadow coming
%		events: the snake on a pole (Num 21:5–9) and Miriam’s being leprous as
%		snow (Num 12:10).
%	}
}

\Verse{10}
{
	\hE{וַיֹּאמֶר מֹשֶׁה אֶל יְהֹוָה בִּי אֲדֹנָי לֹא אִישׁ דְּבָרִים אָנֹכִי גַּם מִתְּמוֹל גַּם מִשִּׁלְשֹׁם גַּם מֵאָז דַּבֶּרְךָ אֶל עַבְדֶּךָ כִּי כְבַד פֶּה וּכְבַד לָשׁוֹן אָנֹכִי׃}
}
{
	\eE{
		And Moses said to YHWH, “Please, my Lord, I’m not a man of
		words. Neither yesterday nor the day before—nor since you
		spoke to your servant! Because I’m heavy of mouth and
		heavy of tongue.”
	}
}
{
%	\footnote{
%		I’m not a man of words. These words, Moses’ fourth response, will
%		reverberate at the beginning of Moses’ last speech. (See the comment on
%		Deut 1:1.) Footnote 4:10. heavy of mouth and heavy of tongue. This is
%		frequently taken to mean some sort of speech defect. It is frequently
%		translated as “slow of speech”; and a famous midrash goes so far as to
%		recount the origin of Moses’ speech impediment in a story that pictures
%		him burning his tongue in his infancy in the Egyptian palace. More
%		cautiously, we are best advised to seek the meaning of a biblical idiom
%		by, first of all, observing where else it occurs in Scripture. “Heavy of
%		tongue” occurs in one other place, Ezek 3:5–7. There YHWH tells Ezekiel
%		that he is not being sent to peoples who are “deep of lip and heavy of
%		tongue,” whose words Ezekiel cannot understand. YHWH says, ironically,
%		that such peoples would listen, but the house of Israel will not listen!
%		In that context, “heavy of tongue” refers to nations who speak foreign
%		languages. It has therefore been suggested that Moses’ protest here in
%		Exodus is that he does not speak Egyptian. This is difficult to defend,
%		though, given the explicit report in Exodus 2 that Moses has been raised
%		in the Egyptian court. Still, the meaning of “heavy of tongue” as
%		referring to speaking a foreign language fits our context in Exodus 4, I
%		believe, because YHWH has told Moses to gather and speak to the elders
%		of Israel. Moses’ protest may perhaps be best understood, then, as being
%		on the grounds that he does not yet speak Hebrew! God’s response in fact
%		confirms that the problem for Moses is speaking “to the people,” not to
%		the Egyptians (4:16). And the final confirmation is that in the first
%		meeting with the people’s elders “Aaron spoke all the words that YHWH
%		had spoken to Moses” (4:30); but in the first meeting with Pharaoh, both
%		Moses and Aaron speak (5:1,3). Footnote 4:10. nor since you spoke to
%		your servant. Another understanding of “I’m not a man of words” and
%		“heavy of mouth and heavy of tongue” is that Moses is saying that he is
%		not eloquent. And this is a wonderful irony, because his humble wording
%		here is in fact an eloquent formulation (“Neither yesterday nor the day
%		before—nor since you spoke to your servant!”). Moses, as it were,
%		undermines his own denial. Whether “heavy of mouth and heavy of tongue”
%		means a speech defect, a foreign language, or a lack of eloquence, a
%		question remains: why choose a man who has a speaking problem for a task
%		that requires speaking?! The lesson may be (as with Jacob) that God can
%		work through anyone. Or it may be that it takes divine insight to
%		comprehend what would make someone the best person for a task. Perhaps
%		it is precisely Moses’ heaviness of mouth that makes his acts most
%		impressive in Egypt, rather than their being performed by a person who
%		heralds each of the miraculous plagues with a grand speech. And perhaps
%		such a greater speaker would shift the focus too much from God to the
%		human himself. Also, one can learn from the divine choice that one
%		should not judge a person’s ability to do a task too hastily or by the
%		most obvious characteristics. Even an individual who has a disadvantage
%		may be the best suited to be successful.
%	}
}

\Verse{11}
{
	\hE{וַיֹּאמֶר יְהֹוָה אֵלָיו מִי שָׂם פֶּה לָאָדָם אוֹ מִי יָשׂוּם אִלֵּם אוֹ חֵרֵשׁ אוֹ פִקֵּחַ אוֹ עִוֵּר הֲלֹא אָנֹכִי יְהֹוָה׃}
}
{
	\eE{
		And YHWH said to him, “Who set a mouth for humans? Or who
		will set a mute or deaf or seeing or blind? Is it not I,
		YHWH?
	}
}
{
}

\Verse{12}
{
	\hE{וְעַתָּה לֵךְ וְאָנֹכִי אֶהְיֶה עִם פִּיךָ וְהוֹרֵיתִיךָ אֲשֶׁר תְּדַבֵּר׃}
}
{
	\eE{
		And now go, and I’ll be with your mouth, and I’ll instruct
		you what you shall speak.”
	}
}
{
%	\footnote{
%		I’ll be. Hebrew ’HYH. Note the significance of the first recurrence of
%		the word here. It also came in 3:12, before the revelation of the divine
%		name. It will come again in 4:15. So the answer to Moses’ first, second,
%		third, and fourth responses is, in essence, the same: ’HYH. I Am. I’ll
%		Be with you.
%	}
}

\Verse{13}
{
	\hE{וַיֹּאמֶר בִּי אֲדֹנָי שְׁלַח נָא בְּיַד תִּשְׁלָח׃}
}
{
	\eE{
		And he said, “Please, my Lord, send by the hand you’ll
		send.”
	}
}
{
%	\footnote{
%		send by the hand you’ll send. There are two opposite meanings that this
%		strange clause may have. It may mean: send whoever—which is to say,
%		anyone but me. Or it may mean that Moses now acquiesces to the divine
%		will, in the sense of “Very well, then, send the one you want.” The fact
%		that God is angry at Moses in the next verse fits more naturally with
%		the former meaning. In either case, the one common element of the two
%		meanings is: they both denote that Moses is out of excuses. God has
%		nullified each of his four attempts to avoid this task. He can only
%		acquiesce or beg off. He cannot raise problems or objections. This is
%		also further evidence that fear of the task is what has motivated Moses
%		all along. He is forced now to the essential point. It is not “Who am I”
%		or “What shall I tell them?” or “They won’t believe me.” It is Moses’
%		reluctance about the hand whom God has chosen to send. Through all five
%		responses, Moses’ reluctance stands out and perhaps comes as a surprise
%		to many who find it difficult to imagine someone’s receiving a prophetic
%		commission directly from God and trying to avoid it. Yet this image of
%		the reluctant prophet recurs in several other places in the Bible
%		(Elijah, Jonah, Jeremiah). It does not appear to me to be simply a
%		recurring literary motif in the way that some understand, for example,
%		the “wise courtier” theme. It seems rather to reflect a shared
%		conviction, by a number of the biblical authors, that it is not
%		necessarily an honor or a joy to be a prophet. It is a burden, almost
%		beyond human endurance. Here, in the book that introduces the role of
%		the prophet as the individual who brings a revelation to a community, we
%		learn from the beginning that this is not attractive. To hear the voice
%		of God is frightening. To experience the divine power is terrifying. To
%		deliver a divine message to the community is thankless, frustrating, and
%		occasionally even dangerous.
%	}
}

\Verse{14}
{
	\hE{וַיִּחַר אַף יְהֹוָה בְּמֹשֶׁה וַיֹּאמֶר הֲלֹא אַהֲרֹן אָחִיךָ הַלֵּוִי יָדַעְתִּי כִּי דַבֵּר יְדַבֵּר הוּא וְגַם הִנֵּה הוּא יֹצֵא לִקְרָאתֶךָ וְרָאֲךָ וְשָׂמַח בְּלִבּוֹ׃}
}
{
	\eE{
		And YHWH’s anger flared at Moses, and He said, “Isn’t
		Aaron your Levite brother? I knew that he will speak! And
		also here he is, coming out toward you! And he’ll see you
		and be happy in his heart.
	}
}
{
%	\footnote{
%		your Levite brother. This means a fellow Levite. The text does not yet
%		identify Aaron as Moses’ actual brother. That is not stated explicitly
%		until Exod 6:20. After that, Aaron is referred to as “your brother”
%		(7:1,2) and never again as “your Levite brother.” Footnote 4:14. I knew
%		that he will speak! This has usually been taken to mean only that God is
%		aware that Aaron is a capable speaker. But the tenses and context
%		possibly indicate that it is much stronger: that God is informing Moses
%		that He knew all along that Aaron would have to do the talking for
%		Moses. This would only be a case of divine foreknowledge (knowing
%		something in advance), not of God preordaining or manipulating human
%		affairs. This in turn opens the question of why God would be angry at
%		Moses’ protest. The ambiguity here is tied to the complexity of Hebrew
%		verb tenses (which are not simply past, present, and future, as they are
%		often taught), and so the meaning of this verse remains uncertain.
%	}
}

\Verse{15}
{
	\hE{וְדִבַּרְתָּ אֵלָיו וְשַׂמְתָּ אֶת הַדְּבָרִים בְּפִיו וְאָנֹכִי אֶהְיֶה עִם פִּיךָ וְעִם פִּיהוּ וְהוֹרֵיתִי אֶתְכֶם אֵת אֲשֶׁר תַּעֲשׂוּן׃}
}
{
	\eE{
		And you’ll speak to him and set the words in his mouth,
		and I, I shall be with your mouth and with his mouth, and
		I shall instruct you what you shall do.
	}
}
{
}

\Verse{16}
{
	\hE{וְדִבֶּר הוּא לְךָ אֶל הָעָם וְהָיָה הוּא יִהְיֶה לְּךָ לְפֶה וְאַתָּה תִּהְיֶה לּוֹ לֵאלֹהִים׃}
}
{
	\eE{
		And he will speak for you to the people. And it will be:
		he will become a mouth for you, and you will become a god
		for him!
	}
}
{
%	\footnote{
%		you will become a god for him. No matter how figuratively we take this
%		image, it is extraordinary. To speak of a human, Moses, in terms of the
%		divine is awesome, and it will recur later when God tells Moses, “I’ve
%		made you a god to Pharaoh!” (7:1). The idea that Moses will in any way
%		be godlike in other humans’ perceptions sets up a theme that will
%		continue after the exodus: Moses must repeatedly insist to the people
%		that their complaints are against God and not against him. This in turn
%		recalls Joseph’s insistence to the wine steward, the baker, and Pharaoh
%		that “Not I. God will answer . . . ” (Gen 41:16).
%	}
}

\Verse{17}
{
	\hE{וְאֶת הַמַּטֶּה הַזֶּה תִּקַּח בְּיָדֶךָ אֲשֶׁר תַּעֲשֶׂה בּוֹ אֶת הָאֹתֹת׃}
}
{
	\eE{
		And you shall take this staff in your hand, by which
		you’ll do the signs.”
	}
}
{
}

\Verse{18}
{
	\hE{וַיֵּלֶךְ מֹשֶׁה וַיָּשׁב אֶל יֶתֶר חֹתְנוֹ וַיֹּאמֶר לוֹ אֵלְכָה נָּא וְאָשׁוּבָה אֶל אַחַי אֲשֶׁר בְּמִצְרַיִם וְאֶרְאֶה הַעוֹדָם חַיִּים וַיֹּאמֶר יִתְרוֹ לְמֹשֶׁה לֵךְ לְשָׁלוֹם׃}
}
{
	\eE{
		And Moses went. And he went back to Jether, his
		father-in-law, and said to him, “Let me go so I may go
		back to my brothers who are in Egypt and see if they’re
		still living.” And Jethro said to Moses, “Go in peace.”
	}
}
{
%	\footnote{
%		Jether. His name is Jethro. The change to yeter here may be part of a
%		wordplay to come later (10:5; see the comment there).
%	}
}

\Verse{19}
{
	\hJ{וַיֹּאמֶר יְהֹוָה אֶל מֹשֶׁה בְּמִדְיָן לֵךְ שֻׁב מִצְרָיִם כִּי מֵתוּ כּל הָאֲנָשִׁים הַמְבַקְשִׁים אֶת נַפְשֶׁךָ׃}
}
{
	\eJ{
		And YHWH said to Moses in Midian, “Go. Go back to Egypt,
		because all the people who sought your life have died.”
	}
}
{
}

\Verse{20}
{
	\hJ{וַיִּקַּח מֹשֶׁה אֶת אִשְׁתּוֹ וְאֶת בָּנָיו וַיַּרְכִּבֵם עַל הַחֲמֹר וַיָּשׁב אַרְצָה מִצְרָיִם}
	\hE{וַיִּקַּח מֹשֶׁה אֶת מַטֵּה הָאֱלֹהִים בְּיָדוֹ׃}
}
{
	\eJ{And Moses took his wife and his sons and rode them on an
		ass, and he went back to the land of Egypt.}
	\eE{And Moses took
		the staff of God in his hand.}
}
{
}

\Verse{21}
{
	\hE{וַיֹּאמֶר יְהֹוָה אֶל מֹשֶׁה בְּלֶכְתְּךָ לָשׁוּב מִצְרַיְמָה רְאֵה כּל הַמֹּפְתִים אֲשֶׁר שַׂמְתִּי בְיָדֶךָ וַעֲשִׂיתָם לִפְנֵי פַרְעֹה}
	\hR{וַאֲנִי אֲחַזֵּק אֶת לִבּוֹ וְלֹא יְשַׁלַּח אֶת הָעָם׃}
}
{
	\eE{And YHWH said to Moses, “When you’re going to go back to
		Egypt, see all the wonders that I’ve set in your hand, and
		you shall do them in front of Pharaoh.}
	\eR{And I: I’ll
		strengthen his heart, and he won’t let the people go.}
}
{
}

\Verse{22}
{
	\hE{וְאָמַרְתָּ אֶל פַּרְעֹה כֹּה אָמַר יְהֹוָה בְּנִי בְכֹרִי יִשְׂרָאֵל׃}
}
{
	\eE{
		And you shall say to Pharaoh, ‘YHWH said this: My child,
		my firstborn, is Israel.
	}
}
{
%	\footnote{
%		My child, my firstborn, is Israel. God tells Moses to say this to
%		Pharaoh, but Moses is never reported as saying it. He also is never
%		reported as saying the rest of this message: “Should you refuse to let
%		it go, here: I’m killing your child, your firstborn!” Why does he not
%		say this to Pharaoh? It is apparently because his first words to Pharaoh
%		do not work. Moses begins by speaking strongly to Pharaoh: “YHWH said,
%		‘Let my people go.’” But Pharaoh, unmoved, replies, “I don’t know YHWH,
%		and I won’t let Israel go.” Moses then changes his approach
%		dramatically. He starts over and this time does not use the commanding
%		wording (the imperative), but rather the polite particle of request
%		(Hebrew n’) and the gentler form: “please that we might go” (the
%		cohortative). Moses has backed off somewhat and taken a more
%		conciliatory approach rather than taking the even tougher approach of
%		threatening Pharaoh with the death of his own son. This conciliatory
%		wording is given to Moses at the burning bush as well (3:18), so it
%		appears that God is giving Moses both the tough and the gentle options.
%		In any negotiation—in law, business, or any human relationship,
%		including marriage—one must use wisdom to know when to exercise each
%		approach.
%	}
}

\Verse{23}
{
	\hE{וָאֹמַר אֵלֶיךָ שַׁלַּח אֶת בְּנִי וְיַעַבְדֵנִי וַתְּמָאֵן לְשַׁלְּחוֹ הִנֵּה אָנֹכִי הֹרֵג אֶת בִּנְךָ בְּכֹרֶךָ׃}
}
{
	\eE{
		And I’ve said to you: Let my child go and serve me. And
		should you refuse to let it go, here: I’m killing your
		child, your firstborn!’”
	}
}
{
}

\Verse{24}
{
	\hJ{וַיְהִי בַדֶּרֶךְ בַּמָּלוֹן וַיִּפְגְּשֵׁהוּ יְהֹוָה וַיְבַקֵּשׁ הֲמִיתוֹ׃}
}
{
	\eJ{
		And he was on the way, at a lodging place, and YHWH met
		him, and he asked to kill him.
	}
}
{
%	\footnote{
%		he asked to kill him. No one knows what the episode at the lodging place
%		means. The explanation that accounts for more of the connections than
%		any other is that of my colleague William Propp: The reason that God
%		seeks to put Moses to death is because Moses still bears the bloodguilt
%		for having killed an Egyptian. The blood of the circumcision serves to
%		expiate the guilt, making it possible for Moses to live and return to
%		Egypt. The word for “blood” in this story is the plural dmîm, which is
%		usually associated with bloodguilt. The word for “bridegroom” also means
%		“circumcision” in Arabic (and in some cultures circumcision is performed
%		just before marriage). Propp suggests that the original
%		bridegroom/circumcision connection had been lost by the time this story
%		was composed, and so it was created to explain an old expression: tan
%		dmîm, “bridegroom of blood.” Also, this use of blood to avert death is a
%		foreshadowing of the tenth plague in Egypt. Still, like all explanations
%		I have seen, this leaves the question of why God would choose Moses,
%		make a miraculous appearance to him, insist on his going to Egypt
%		despite his five attempts to avoid it, and then seek to kill him. I
%		therefore understand the confused use of the pronoun “he” in this story
%		in a different way. All commentators recognize that it is unclear when
%		it refers to God, to Moses, or to Moses’ son. But they all take the
%		first “he” to refer to God: He (God) sought to kill Moses. I am raising
%		the possibility (reflected in my translation: “he asked to kill him”)
%		that it means that Moses is asking God to take his life (rather than
%		send him to Egypt). This is consistent with another time in Moses’ life
%		when he will ask God to kill him (Num 11:15, where he says emphatically:
%		“Kill me!”). And it fits with a model of prophets who ask God to take
%		their lives—Elijah (1 Kings 19:4) and Jonah (4:4)—or say that they
%		prefer death to being prophets—Jeremiah (20:14–18). (Note that these are
%		the same three prophets who, like Moses, are famous for being reluctant;
%		see the comment on 4:13.) If Moses is seeking his own death here, then
%		Zipporah’s action might be understood as her passionate response to her
%		husband’s wish to die: You have a wife! You have a son, who should live
%		to marry and be a part of the covenant!
%	}
}

\Verse{25}
{
	\hJ{וַתִּקַּח צִפֹּרָה צֹר וַתִּכְרֹת אֶת ערְלַת בְּנָהּ וַתַּגַּע לְרַגְלָיו וַתֹּאמֶר כִּי חֲתַן דָּמִים אַתָּה לִי׃}
}
{
	\eJ{
		And Zipporah took a flint and cut her son’s foreskin and
		touched his feet, and she said, “Because you’re a
		bridegroom of blood to me.”
	}
}
{
}

\Verse{26}
{
	\hJ{וַיִּרֶף מִמֶּנּוּ אָז אָמְרָה חֲתַן דָּמִים לַמּוּלֹת׃}
}
{
	\eJ{
		And he held back from Him. Then she said, “A bridegroom of
		blood for circumcisions.”
	}
}
{
%	\footnote{
%		And he held back from Him. Meaning that Moses stopped asking God to kill
%		him.
%	}
}

\Verse{27}
{
	\hE{וַיֹּאמֶר יְהֹוָה אֶל אַהֲרֹן לֵךְ לִקְרַאת מֹשֶׁה הַמִּדְבָּרָה וַיֵּלֶךְ וַיִּפְגְּשֵׁהוּ בְּהַר הָאֱלֹהִים וַיִּשַּׁק לוֹ׃}
}
{
	\eE{
		And YHWH said to Aaron, “Go toward Moses, to the
		wilderness.” And he went, and he met him in the Mountain
		of God, and he kissed him.
	}
}
{
}

\Verse{28}
{
	\hE{וַיַּגֵּד מֹשֶׁה לְאַהֲרֹן אֵת כּל דִּבְרֵי יְהֹוָה אֲשֶׁר שְׁלָחוֹ וְאֵת כּל הָאֹתֹת אֲשֶׁר צִוָּהוּ׃}
}
{
	\eE{
		And Moses told Aaron all YHWH’s words that He had sent him
		and all the signs that He had commanded him.
	}
}
{
}

\Verse{29}
{
	\hE{וַיֵּלֶךְ מֹשֶׁה וְאַהֲרֹן וַיַּאַסְפוּ אֶת כּל זִקְנֵי בְּנֵי יִשְׂרָאֵל׃}
}
{
	\eE{
		And Moses went, and Aaron, and they gathered all the
		elders of the children of Israel,
	}
}
{
}

\Verse{30}
{
	\hE{וַיְדַבֵּר אַהֲרֹן אֵת כּל הַדְּבָרִים אֲשֶׁר דִּבֶּר יְהֹוָה אֶל מֹשֶׁה וַיַּעַשׂ הָאֹתֹת לְעֵינֵי הָעָם׃}
}
{
	\eE{
		and Aaron spoke all the words that YHWH had spoken to
		Moses. And he did the signs before the people’s eyes.
	}
}
{
}

\Verse{31}
{
	\hE{וַיַּאֲמֵן הָעָם וַיִּשְׁמְעוּ כִּי פָקַד יְהֹוָה אֶת בְּנֵי יִשְׂרָאֵל וְכִי רָאָה אֶת ענְיָם וַיִּקְּדוּ וַיִּשְׁתַּחֲווּ׃}
}
{
	\eE{
		And the people believed, and they heard that YHWH had
		taken account of the children of Israel and that He had
		seen their degradation. And they knelt and bowed.
	}
}
{
}
